%%% FIELD proof-total-certificate
Fix effective \(K\)-valued matrices \(B,Q\), the supplied graph \(G\), and the known integer \(m\geq2\).  For a compatible structural pair \((A,D)\), the causal target is the fine-grid transfer \(h_{uv}(z)=A_{vu}z/(1-A_{vv}z)\).  The set \(\mathcal H_{uv}(B,Q;G)\), by contrast, is the image of the compatible structural fiber under this transfer map.  We prove that the certificate decides whether this image is empty, a singleton, or contains two unequal attained transfers.

Define the finite coordinate set
\[
 \mathcal I_G
 =\{(i,i):i\in V\}\cup
   \{(i,j):j<i\text{ and }j\to i\text{ in }G\}.
\]
For every \((i,j)\in\mathcal I_G\), introduce a real variable \(a_{ij}\), and introduce real variables \(d_1,\ldots,d_d\).  Let \(\widetilde A(a)\) have entry \(a_{ij}\) on \(\mathcal I_G\) and zero elsewhere, and put \(\widetilde D(d)=\operatorname{diag}(d_1,\ldots,d_d)\).  Thus lower triangularity and graph support are imposed by the coordinate parametrization itself.

Consider the finite system
\[
 [\widetilde A(a)^m]_{ij}=B_{ij}\qquad(i,j\in V),                 \tag{1}
\]
\[
 \left[\sum_{r=0}^{m-1}\widetilde A(a)^r\widetilde D(d)
                    (\widetilde A(a)^r)^\top\right]_{ij}
 =Q_{ij}\qquad(i,j\in V),                                        \tag{2}
\]
and
\[
 -1<a_{ii}<1\qquad(i\in V),
 \qquad d_i>0\qquad(i\in V).                                    \tag{3}
\]
Every left side in (1)--(2) is a polynomial in finitely many variables with coefficients in \(K\), by \(\mathrm{ass{:}effective\mbox{-}input}\).  Because \(\widetilde A(a)\) is triangular,
\[
 \det(\lambda I-\widetilde A(a))=\prod_{i=1}^d(\lambda-a_{ii}),
 \qquad \rho(\widetilde A(a))=\max_i|a_{ii}|.
\]
Hence the first inequalities in (3) are equivalent to stability.  Moreover, for every nonzero \(x\in\mathbb R^d\),
\[
 x^\top\widetilde D(d)x=\sum_{i=1}^d d_i x_i^2>0
\]
if and only if every \(d_i>0\); thus the second inequalities in (3) are equivalent to \(\widetilde D(d)\succ0\).  It follows directly from \(\mathrm{def{:}compatible\mbox{-}fiber}\) that
\[
 (a,d)\longmapsto (\widetilde A(a),\widetilde D(d))               \tag{4}
\]
is a bijection from the realization of (1)--(3) onto \(\mathcal S(B,Q;G)\).  In particular, this encoding preserves the strict stability and positivity conditions and introduces no boundary points.

The diagonal equations in (1) are
\[
 a_{ii}^{m}=B_{ii}\qquad(i\in V).                                \tag{5}
\]
The effective arithmetic, sign tests, and real-root isolation supplied by the definition of the effective real closed field \(K\) enumerate their real roots and decide (3).  For odd \(m\), \(x\mapsto x^m\) is strictly increasing and onto \(\mathbb R\), so (5) has exactly one real root for each \(i\).  For even \(m\), (5) has no real root when \(B_{ii}<0\), the single root zero when \(B_{ii}=0\), and the two roots \(\pm B_{ii}^{1/m}\) when \(B_{ii}>0\).  There are therefore finitely many surviving stable diagonal branches, all exactly represented in \(K\).  If none survives, (4) proves \(\mathcal S(B,Q;G)=\varnothing\).

Fix a surviving branch and process off-diagonal pairs \((v,u)\) in increasing order of \(v-u\), as in \(\mathrm{def{:}graph\mbox{-}recursion}\).  By \(\mathrm{lem{:}path\mbox{-}expansion}\),
\[
 a_{vu}C_m(a_{vv},a_{uu})
 =B_{vu}-R_{vu,m}(\widetilde A(a)).                               \tag{6}
\]
Every off-diagonal coordinate on the right of (6) has smaller endpoint gap.  The diagonal entries are fixed on the branch, so
\[
 c_{vu}:=C_m(a_{vv},a_{uu})\in K.                                \tag{7}
\]
The zero-set characterization \(\mathrm{lem{:}divided\mbox{-}difference\mbox{-}zero\mbox{-}set}\), together with exact sign tests in \(K\), decides whether \(c_{vu}=0\).  Thus the odd case \(a_{vv}=a_{uu}=0\) and the even case \(a_{vv}=-a_{uu}\) are retained as singular cases and are never divided through.

If \(u\to v\) is allowed by \(G\) and \(c_{vu}\ne0\), equation (6) is equivalent to the unique substitution
\[
 a_{vu}
 =\frac{B_{vu}-R_{vu,m}(\widetilde A(a))}{c_{vu}}.                \tag{8}
\]
Division is legitimate because of the tested condition \(c_{vu}\ne0\).  Since the inverse of this fixed nonzero element belongs to the field \(K\), substitution into later transition equations and into (2) leaves polynomial equations over \(K\) in the residual variables.  If \(c_{vu}=0\), the coordinate is not eliminated and the induced equation
\[
 B_{vu}-R_{vu,m}(\widetilde A(a))=0                               \tag{9}
\]
is retained.  If \(u\to v\) is forbidden, the parametrization fixes \(a_{vu}=0\), and its transition compatibility equation, which is again (9), is retained regardless of \(c_{vu}\).

We now prove exactness of the recursion by induction on the endpoint gap.  At an allowed nonsingular coordinate, (6) shows that every solution before substitution has the value (8), and direct replacement shows that every solution of the substituted system lifts uniquely through (8).  At a singular or forbidden coordinate, no residual variable or compatibility equation is discarded.  The induction is valid because \(\mathrm{lem{:}path\mbox{-}expansion}\) places only smaller-gap off-diagonal coordinates in the remainder.  After the finitely many pairs have been processed, every transition equation not used for an exact substitution, every covariance equation in (2), and every strict inequality in (3) remains imposed.  If \(\mathfrak R_b\) denotes the residual semialgebraic realization on branch \(b\), and \(\operatorname{Rec}_b\) denotes reconstruction by all substitutions, then
\[
 \mathcal S(B,Q;G)
 =\bigcup_b\operatorname{Rec}_b(\mathfrak R_b).                  \tag{10}
\]
This is literal set equality, not equality after algebraic or topological closure.

Each \(\mathfrak R_b\) is a finite Boolean combination of polynomial equalities and strict polynomial inequalities over \(K\).  There are finitely many branches and variables.  By \(\mathrm{lem{:}effective\mbox{-}cad\mbox{-}decision\mbox{-}sampling}\), CAD terminates on every branch, decides whether \(\mathfrak R_b\) is empty, and returns an exact represented sample point when it is nonempty, including for open or positive-dimensional realizations.  Equation (10) therefore implies that all branch systems are empty if and only if \(\mathcal S(B,Q;G)=\varnothing\).  This proves termination, soundness, and completeness of the incompatibility output.

It remains to classify the transfer fiber.  Take two compatible points and write
\[
 p=A_{vu},\quad q=A_{vv},\qquad p'=A'_{vu},\quad q'=A'_{vv}.
\]
Stability gives \(|q|<1\) and \(|q'|<1\).  Therefore, for every \(z\in\mathbb T\), both denominators \(1-qz\) and \(1-q'z\) are nonzero, and \(z\ne0\).  Cross multiplication is consequently valid, and the two transfers agree on \(\mathbb T\) if and only if
\[
 p(1-q'z)=p'(1-qz)\qquad\text{for every }z\in\mathbb T.          \tag{11}
\]
The difference between the two sides is
\[
 (p-p')+(p'q-pq')z.                                               \tag{12}
\]
Evaluating (12) at \(z=1\) and \(z=-1\) shows that it vanishes on \(\mathbb T\) if and only if both coefficients vanish.  Hence
\[
 h_{p,q}=h_{p',q'}
 \quad\Longleftrightarrow\quad
 p=p'\ \text{ and }\ pq'=p'q,                                   \tag{13}
\]
and, because the coefficients are real,
\[
 h_{p,q}\ne h_{p',q'}
 \quad\Longleftrightarrow\quad
 \Delta(p,q,p',q')
 :=(p-p')^2+(pq'-p'q)^2>0.                                       \tag{14}
\]
In particular, (13) correctly recognizes that every pair \((0,q)\) represents the same zero transfer; equality of coefficient pairs is not required.

For every ordered pair of diagonal branches, form two copies of their complete residual systems and append the strict polynomial inequality (14), after reconstruction.  By (10) and (14), the finite union of these two-copy systems is feasible if and only if two points of \(\mathcal S(B,Q;G)\) have unequal transfers.  The cited CAD lemma decides feasibility and, when feasible, supplies exact residual sample points.  Reconstruction then gives attained pairs
\[
 (A,D),(A',D')\in\mathcal S(B,Q;G)
\]
such that
\[
 (A_{vu}-A'_{vu})^2
 +(A_{vu}A'_{vv}-A'_{vu}A_{vv})^2>0.
\]
Their strict stability and innovation positivity follow from their one-copy systems.  This proves both soundness and completeness of the ambiguity output.

If every two-copy system is infeasible, (14) says that all compatible points yield the same transfer.  An exact sample from any nonempty one-copy branch supplies \((p,q)\), so the certificate returns the common exact rational function
\[
 z\longmapsto\frac{pz}{1-qz}.
\]
Conversely, if \(\mathcal H_{uv}(B,Q;G)\) is a singleton, no two-copy system can satisfy (14), and the certificate returns that common function.  Therefore the singleton output is equivalent to point identification of the causal transfer, and the three outputs in the theorem are exhaustive.

The preceding equivalences also cover all listed exceptional cases.  Zero diagonal roots remain among (5) and enter a residual singular system whenever their divided difference vanishes.  A forbidden queried edge, or an allowed queried edge forced to zero by compatibility, is handled by (13).  Strict inequalities and the cited CAD result cover open and positive-dimensional fibers.  Duplicate branches or branch mergers do not affect the conclusion because (10) is a union and the two-copy test ranges over every ordered pair of branches.  No rank, overlap, conditioning, boundary passage, or limit argument is used; the only divisions are the explicitly certified divisions in (8).

Finally, the two prescribed exact stress tests confirm that the construction does not lose singular fibers or the even-factor sign pair.  For the first, take \(d=2\), \(m=3\), let \(G\) contain \(1\to2\), and set \(B=0\) and \(Q=I_2\).  Define
\[
 A_{\pm}=\begin{pmatrix}0&0\\ \pm\tfrac12&0\end{pmatrix},
 \qquad
 D_{\pm}=\operatorname{diag}\!\left(1,\tfrac34\right).
\]
Then \(A_{\pm}^2=0\), and hence
\[
 A_{\pm}^3=0=B,
 \qquad
 \sum_{r=0}^{2}A_{\pm}^rD_{\pm}(A_{\pm}^r)^\top
 =D_{\pm}+A_{\pm}D_{\pm}A_{\pm}^\top
 =\operatorname{diag}\!\left(1,\tfrac34\right)
  +\operatorname{diag}\!\left(0,\tfrac14\right)
 =I_2.
\]
Both pairs are strictly interior because \(\rho(A_{\pm})=0<1\) and \(D_{\pm}\succ0\), while their transfers are \(z/2\) and \(-z/2\).  Thus the exact two-copy test must retain two attained, interior, transfer-distinct witnesses.

For the second test, take \(d=2\), \(m=2\), and any compatible \((A,D)\) with \(A_{21}\ne0\).  The map \(A\mapsto-A\) preserves triangular and graph-support zeros, stability, and positivity of \(D\), while
\[
 (-A)^2=A^2,
 \qquad
 D+(-A)D(-A)^\top=D+ADA^\top.
\]
Thus \((-A,D)\) is compatible with the same \((B,Q)\).  With \(p=A_{21}\), \(q=A_{22}\), \(p'=-p\), and \(q'=-q\),
\[
 \Delta(p,q,-p,-q)
 =(2p)^2+\{p(-q)-(-p)q\}^2
 =4p^2>0.
\]
Every nonzero-edge two-variable \(m=2\) compatible case therefore retains its global sign pair.
